\documentclass[11pt,a4paper]{article}

\usepackage[T1]{fontenc}
\usepackage[utf8]{inputenc}
\usepackage{lmodern}
\usepackage{textcomp}
\usepackage{microtype}
\usepackage[a4paper,margin=2.3cm]{geometry}
\usepackage{xcolor}
\usepackage{amsmath}
\usepackage{booktabs,longtable,tabularx,array}
\usepackage{enumitem}
\usepackage{hyperref}
\usepackage{listings}
\usepackage{fancyhdr}
\usepackage{titlesec}
\usepackage{parskip}
\usepackage{pifont}

\definecolor{meopblue}{HTML}{143C6B}
\definecolor{meoplight}{HTML}{EAF1F8}
\definecolor{meopmid}{HTML}{64748B}
\definecolor{meopgreen}{HTML}{2E7D32}
\definecolor{meopamber}{HTML}{B26A00}
\definecolor{meopred}{HTML}{A3312B}
\definecolor{codebg}{HTML}{F7F8FA}

\hypersetup{
  colorlinks=true,
  linkcolor=meopblue,
  urlcolor=meopblue,
  citecolor=meopblue,
  pdftitle={MEOP temperature and salinity comparison and adjustment framework},
  pdfauthor={MEOP project}
}

\titleformat{\section}{\Large\bfseries\color{meopblue}}{\thesection}{0.7em}{}
\titleformat{\subsection}{\large\bfseries\color{meopblue}}{\thesubsection}{0.7em}{}
\titleformat{\subsubsection}{\normalsize\bfseries\color{meopblue}}{\thesubsubsection}{0.7em}{}

\pagestyle{fancy}
\fancyhf{}
\fancyhead[L]{\textcolor{meopmid}{MEOP T/S comparison and adjustment}}
\fancyhead[R]{\textcolor{meopmid}{August 22, 2026}}
\fancyfoot[C]{\textcolor{meopmid}{\thepage}}
\setlength{\headheight}{15pt}

\lstdefinestyle{meopcode}{
  basicstyle=\ttfamily\small,
  backgroundcolor=\color{codebg},
  frame=single,
  rulecolor=\color{black!10},
  framerule=0.6pt,
  columns=fullflexible,
  keepspaces=true,
  showstringspaces=false,
  breaklines=true,
  upquote=true
}

\newcommand{\done}{\textcolor{meopgreen}{\ding{51}}}
\newcommand{\partdone}{\textcolor{meopamber}{\ding{108}}}
\newcommand{\missing}{\textcolor{meopred}{\ding{55}}}
\newcommand{\pkg}[1]{\texttt{\detokenize{#1}}}
\newcommand{\goodbox}[1]{%
  \noindent\colorbox{meoplight}{\parbox{0.955\textwidth}{\vspace{0.2cm}#1\vspace{0.2cm}}}}

\begin{document}

\begin{titlepage}
  \centering
  \vspace*{1.4cm}
  {\Huge\bfseries\color{meopblue}From comparison to defensible\\[0.2em]
  temperature and salinity adjustments\par}
  \vspace{0.7cm}
  {\Large A reusable scientific and software-development framework\par}
  \vspace{1.1cm}

  \goodbox{
    \textbf{Recommended scope.} Build a systematic, profile-by-profile comparison of
    unadjusted MEOP temperature and salinity with independent reference profiles.
    Make the comparison reproducible, quality-controlled, uncertainty-aware, and
    usable at the scale of the MEOP archive. Once that foundation is validated,
    use it to produce \emph{candidate} adjustment coefficients for expert review.
    Automatic modification of the operational coefficient tables is explicitly
    outside the scope of this framework.
  }

  \vspace{1.1cm}
  \begin{tabularx}{0.94\textwidth}{>{\bfseries}p{0.28\textwidth}X}
    Primary objective & A tested space--time matchup and comparison workflow, with documented outputs and a pilot scientific assessment \\
    Optional extension & Conservative T/S adjustment suggestions with uncertainty and explicit accept/review/reject reasons \\
    Existing foundation & Processed MEOP profiles, the CORA 202511 CTD/Argo reference archive, manifest-based acquisition and preparation, explicit batch statuses, coefficient application, and plotting/batch machinery \\
    Main scientific principle & Separate sensor bias from natural ocean variability, and never use an adjustment merely because a regression returns a number \\
    Snapshot date & August 22, 2026 \\
  \end{tabularx}

  \vfill
  {\large MEOP Process toolbox\par}
\end{titlepage}

\pagenumbering{roman}
\tableofcontents
\clearpage
\pagenumbering{arabic}

\section{The decision and the project goal}

\subsection{Recommendation}

The two possible work packages are not truly separate. A method that estimates
temperature or salinity adjustments first needs a scientifically defensible comparison:
which external profiles are relevant, which observations pass quality control, how the
data are brought to a common vertical coordinate, and how much of a difference can be
explained by normal spatial and temporal ocean variability.

The recommended sequence is therefore:

\begin{quote}
\textbf{Primary goal:} develop and validate a systematic comparison framework.\\
\textbf{Gated extension:} use the same framework to estimate candidate adjustments,
but only after the comparison passes agreed scientific tests.
\end{quote}

This scope has a useful minimum result even if adjustment estimation proves too
ambitious: the comparison archive, coverage report, diagnostics and tested matchup code
will remain directly useful for future delayed-mode work. Conversely, starting with a
coefficient estimator without fixing the comparison would risk fitting ocean variability,
reference-data artefacts, or corrections that have already been applied.

\begin{center}
\begin{tabularx}{0.96\textwidth}{>{\bfseries}p{0.22\textwidth}XX}
\toprule
 & Systematic comparison & Adjustment estimation \\
\midrule
Scientific risk & Moderate; assumptions can be exposed and tested & Higher; natural variability can look like instrument drift \\
Standalone value & High even if adjustment estimation is deferred & High only if validation is convincing \\
Existing support & CORA loading, plots and batch runner & Coefficient format and coefficient application \\
Main missing element & Profile-level space--time matchups & Validated estimator, uncertainty and decision rules \\
Recommended role & Core work package & Gated extension after scientific review \\
\bottomrule
\end{tabularx}
\end{center}

\subsection{Scientific questions}

The project should answer the following questions in order:

\begin{enumerate}[leftmargin=1.6em]
  \item For each usable MEOP profile, are there independent reference profiles close
  enough in space and time to make a meaningful comparison?
  \item After applying quality control and interpolating to a common pressure grid,
  what are the robust MEOP-minus-reference temperature and salinity anomalies?
  \item How do these anomalies vary with pressure, time, region, season, matching
  radius and reference selection?
  \item Are the observed patterns compatible with a constant offset, a pressure-dependent
  offset, a time-varying salinity offset, or simply unresolved ocean variability?
  \item When a correction is proposed, does it improve independent held-out comparisons
  and remain stable under reasonable changes to the matching parameters?
\end{enumerate}

\subsection{Definition of success}

The work is successful when it delivers all of the following core items:

\begin{itemize}[leftmargin=1.5em]
  \item a documented and tested CORA inventory/QC reader that preserves profile identity,
  location, time, quality flags and provenance;
  \item an anti-meridian-safe, configurable space--time matchup method;
  \item a stable comparison data product for each tag, containing raw and adjusted
  anomalies, coverage and reference-spread estimates;
  \item reproducible per-tag and multi-tag summaries and plots;
  \item synthetic tests plus a scientifically diverse pilot evaluation;
  \item a batch command that separates usable results, missing or insufficient
  reference data, invalid targets and genuine software failures;
  \item concise user and developer documentation, including known limitations.
\end{itemize}

Candidate adjustment coefficients are a useful extension, but not a condition for a
successful comparison framework. They should only be produced after an explicit
scientific review accepts the comparison results and uncertainty analysis.

\section{Review of what is already in place}

\subsection{Operational processing foundation}

The renovated toolbox is a Python package under \pkg{src/meop_process/}. It already owns
the main MEOP processing path: raw ODV import, location processing, quality control,
low- and high-resolution products, thermal-lag processing, adjustment application,
diagnostics, metadata summaries, publishing helpers and resumable batch execution.
The package exposes installed commands including \pkg{meop-process},
\pkg{meop-compare}, \pkg{meop-compare-batch} and \pkg{meop-publish}.

The repository test suite contains 223 collected tests in this snapshot. Relevant
regressions cover adjustment application, CORA tile loading, calibration-plot
generation, batch resumption and NetCDF comparison. This is a strong software base:
new work should extend it rather than create a disconnected notebook-only method.

\subsection{Local data inventory at the review date}

The working data tree is already large enough for development and retrospective
validation. Counts below describe the local checkout on August 22, 2026; they are a
snapshot rather than a property of the software.

\begin{center}
\begin{tabularx}{0.94\textwidth}{>{\raggedright\arraybackslash}p{0.34\textwidth}X}
\toprule
Item & Count \\
\midrule
Deployment directories under \pkg{data/data_prof/} & 310 \\
Unique tags with processed profile files & 2,717 \\
\pkg{lr0}, \pkg{lr1}, \pkg{hr0}, \pkg{hr1}, or \pkg{hr2} files & 2,717 of each \\
\pkg{fr0} and \pkg{fr1} files & 182 of each \\
Rows in \pkg{table_coeff.csv} & 2,789 \\
Tags represented in \pkg{table_salinity_offsets.csv} & 77 \\
CORA 202511 populated tile files & 540 \\
CORA 202511 profiles & 4,328,083: 2,605,655 Argo and 1,722,428 CTD \\
CORA 202511 coverage & CTD: 1953-04-09 to 2025-06-30; Argo: 1995-09-07 to 2025-06-30 \\
Collected repository tests & 223 \\
\bottomrule
\end{tabularx}
\end{center}

\subsection{Existing pieces directly relevant to the work}

\begin{longtable}{>{\raggedright\arraybackslash}p{0.25\textwidth} >{\centering\arraybackslash}p{0.09\textwidth} >{\raggedright\arraybackslash}p{0.56\textwidth}}
\toprule
\textbf{Component} & \textbf{Status} & \textbf{What it does today} \\
\midrule
\endhead
Processed profile products & \done & One NetCDF file per tag and processing level under \pkg{data/data_prof/<deployment>/}. \\
Runtime reference configuration & \done & \pkg{configs.json} points \pkg{references.cora_dir} to \pkg{CORA_ctd_argo_202511_tiles}. \\
CORA tile reader & \partdone & \pkg{reference/cora.py} inventories actual files, accepts padded and unpadded longitude names, and returns latitude, longitude, JULD, TEMP, PSAL and a shared pressure grid. Calibration loads the union of buffered cells along accepted positions, handles anti-meridian windows circularly, and indexes selected profile rows before materialising TEMP/PSAL. The reader does not yet expose the retained profile provenance required for auditable matchups. \\
CORA source acquisition & \done & \pkg{meop-cora-download} produced a complete, release-pinned CORA 202511 source archive with exact hashed CTD and Argo file lists and a download manifest. \\
CORA tile preparation & \done & \pkg{meop-cora-prepare} applied per-profile type/QC/value rules, normalised time and longitude, interpolated to the 10-dbar grid, retained source/type/branch metadata, and wrote the installed 540-tile archive and preparation manifest. \\
Calibration figures & \partdone & \pkg{plotting/calibration.py} produces T/S context plots, salinity anomaly plots, pressure-band statistics and a linear pressure fit. \\
Single-tag command & \done & \pkg{meop-compare --plot1 <tag>} selects a processed target file and writes calibration PNG/CSV outputs. \\
Batch command & \done & \pkg{meop-compare-batch} discovers tags, resumes version-compatible successes, runs in parallel, logs each tag, distinguishes data-availability outcomes from failures, and writes CSV/Markdown summaries. \\
Regression comparator & \done & \pkg{workflows/compare.py} and directory mode in \pkg{meop-compare} compare two NetCDF outputs for numerical/software regression. This is not an oceanographic matchup method. \\
Coefficient application & \done & \pkg{processing/adjustments.py} reads operator tables, applies T/S corrections and writes adjusted variables, errors and calibration metadata. \\
Coefficient estimation & \partdone & Existing calibration CSVs contain preliminary suggested coefficients, but the scientific comparison used to obtain them is not yet a valid basis for operational decisions. \\
Time-varying salinity adjustment & \done & Existing table values are interpolated by profile index and applied. Their breakpoints are operator supplied; no estimator finds them. \\
Expert-review workflow & \missing & No proposal table, confidence classification, before/after validation report or controlled approval/import step exists. \\
\bottomrule
\end{longtable}

\subsection{What the current calibration workflow actually computes}

For a target tag, the existing command prefers \pkg{hr1}, then \pkg{lr1},
\pkg{hr2}, \pkg{hr0}, \pkg{lr0}, \pkg{fr1} and \pkg{fr0}. Target preflight removes
non-finite/null-island positions and minor spatial components disconnected from a
supported deployment component. It verifies target-only profile and depth support
before opening CORA. The reader then loads only cells intersecting local 5-degree
windows along accepted positions and indexes matching reference rows before
materialising T/S arrays. It projects the retained CORA profiles to the target pressure
grid and computes one CORA median at each pressure level. Every accepted target profile
is compared with this same median profile.

The diagnostic CSV reports robust summaries in fixed pressure bands and an ordinary
linear fit over 200--1000 dbar. If

\begin{equation}
\Delta T(p) = T_{\mathrm{tag}}(p)-T_{\mathrm{CORA,median}}(p)
             \simeq b_T + a_T p,
\end{equation}

the CSV writes candidate values \(T1=1000a_T\) and \(T2=b_T\), with analogous
values for salinity. These are syntactically compatible with the coefficient table,
but the reference median is not matched to each target profile in time or local space.
The numerical compatibility must not be confused with scientific validation.

\subsection{How adjustments are already represented and applied}

The operational table \pkg{data/tables/table_coeff.csv} stores one row per tag with
\pkg{T1}, \pkg{T2}, \pkg{S1}, \pkg{S2}, removal flags and a comment code. The Python
implementation applies

\begin{align}
T_{\mathrm{adjusted}}(p) &= T(p) - 10^{-3}T1\,p - T2, \\
S_{\mathrm{adjusted}}(p,i) &= S(p,i) - 10^{-3}S1\,p - S2 + \delta S(i),
\end{align}

where \(p\) is pressure in dbar, \(i\) is profile index, and \(\delta S(i)\) is an
optional time-varying series interpolated from up to four index/offset pairs in
\pkg{table_salinity_offsets.csv}. The code also updates adjusted-error variables and
NetCDF calibration metadata.

The local coefficient history is valuable validation material: 262 rows have non-zero
\pkg{T2}, 285 have non-zero \pkg{S1}, 1,051 have non-zero \pkg{S2}, and 77 tags have
time-varying salinity-offset records. These are expert decisions, not perfect ground
truth. They can be used to form test strata and to investigate disagreement, but should
not simply be learned and reproduced.

The deprecated MATLAB code retained under \pkg{legacy/toolbox_legacy/} likewise applies
pre-existing coefficients and interpolated salinity offsets. It does not provide the
missing statistically validated estimator. The old external calibration plotting call
is referenced by the legacy wrapper, but its full implementation is not present in this
repository.

\subsection{CORA 202511 reference archive}

The configured reference is the prepared CTD/Argo archive from Copernicus Marine
product \pkg{INSITU_GLO_PHY_TS_DISCRETE_MY_013_001}, dataset version \pkg{202511}.
The download manifest records 34,556 source files: raw CORA \pkg{CT/OC/TE} candidates
for CTD selection and EasyCORA \pkg{PF} files received from Argo DACs. The preparation
manifest records 540 populated 10-degree tiles and 4,328,083 retained profiles:
2,605,655 Argo profiles with \pkg{PROBE_TYPE=5} and 1,722,428 CTD profiles with
\pkg{PROBE_TYPE=2}.

CTD coverage is 1953-04-09 through 2025-06-30; Argo coverage is 1995-09-07 through
2025-06-30. Every tile uses one standard time origin and a 100-level pressure grid from
10 to 1000 dbar in 10-dbar increments. A complete structural audit found no retained
profile with non-finite time or position, no profile outside its named tile, no
unexpected probe type, no malformed pressure grid, no skipped source file, and no
missing QC variable.

The preparation accepts level, position and time QC flags 1 and 2, prefers an adjusted
whole-profile parameter when available, otherwise uses the raw parameter, and requires
both usable temperature and salinity on the target grid. Of 5,271,991 profiles matching
the requested platform types, 400,035 were excluded by the strict time/position policy
and 543,873 lacked sufficient paired T/S support. These exclusions are recorded in the
manifest and should remain part of the scientific provenance.

The acquisition and preparation workflow is reproducible:

\begin{itemize}[leftmargin=1.5em]
  \item \pkg{meop-cora-download plan|download} first saves the complete remote inventory,
  then selects exact global file lists, records size, version and hashes, and downloads
  source files without mixing them with generated tiles;
  \item \pkg{meop-cora-prepare} writes MEOP-compatible 10-degree tiles and a generation
  manifest while retaining \pkg{REFERENCE_KIND}, \pkg{PROBE_TYPE}, platform/reference
  identifiers, source file/class/dataset and the selected raw/adjusted parameter branch.
\end{itemize}

EasyCORA provides best-estimate, vertically subsampled Argo values, which is suitable
for the 10-dbar comparison tiles but must remain visible in provenance. Full commands,
policies and manifest-review checks are in
\pkg{docs/cora_download_and_preparation.md}.

\subsection{Calibration batch status after the reference update}

The calibration products must be regenerated after a reference archive changes. The
current batch command discovers 2,665 tags in enabled deployments, or all 2,717 locally
processed tags with \pkg{--include-disabled}. It writes per-tag context plots, an offset
diagnostic PNG and a diagnostic CSV containing pressure-band summaries and preliminary
\pkg{T1/T2/S1/S2} estimates. Batch method version 4 distinguishes \pkg{no_reference},
\pkg{insufficient_target}, \pkg{insufficient_reference}, \pkg{invalid_target} and
\pkg{failed} from successful results and refuses to reuse incompatible success state.
Success requires, for both
temperature and salinity, at least 10 contributing target profiles, at least three valid
levels in the 400--600 dbar band, finite band and linear estimates, and actual comparison
coverage spanning at least 200--600 dbar. A clean rebuild must remove calibration plots,
offset CSVs and calibration batch state before running
\pkg{meop-compare-batch --include-disabled --force}.

\section{What still needs to be developed}

\subsection{The central missing capability}

The toolbox needs to move from a \emph{track-wide visual background} to an auditable
\emph{profile-level comparison}. Each target profile must have its own reference
neighbourhood, with explicit spatial distance, time separation, reference identity,
quality-control decisions and uncertainty. All downstream summaries and coefficient
suggestions must be traceable to those matchups.

\subsection{Specific gaps found during the review}

\begin{longtable}{>{\raggedright\arraybackslash}p{0.25\textwidth} >{\raggedright\arraybackslash}p{0.34\textwidth} >{\raggedright\arraybackslash}p{0.31\textwidth}}
\toprule
\textbf{Gap} & \textbf{Why it matters} & \textbf{Direction} \\
\midrule
\endhead
No per-profile time/space matching & A broad regional median mixes seasons, fronts and water masses & Match every target profile separately with configurable distance/time limits \\
Adjusted variables are preferred & Existing coefficients can leak into an estimator intended to discover those coefficients & Compare explicit raw \pkg{TEMP}/\pkg{PSAL}; carry adjusted values only as a separate evaluation branch \\
Prepared CORA provenance is not exposed by the calibration loader & Source class, platform identity and raw/adjusted branch cannot yet be audited downstream & Return retained identifiers and branch metadata with the numerical arrays \\
MEOP value and position/time QC are not applied in calibration & Fill-value removal alone is not sufficient & Centralise flag decoding and masks for target and reference data \\
Reference identity is discarded & Self-matches, duplicated profiles and source dependence cannot be audited & Retain platform, cycle, data centre/reference, mode and source file/profile index \\
Reference archive identity is not part of batch invalidation & A successful state can be reused after the reference data change & Include source/preparation manifest hashes in the method configuration and batch state \\
Reference time is not used by the calibration method & A track-wide median mixes observations from different seasons and years & Decode the normalised \pkg{JULD} coordinate, match in time and test known dates \\
One value per pressure level is treated as independent & Confidence becomes far too optimistic because levels within a profile are correlated & Bootstrap or split by profile/time block, not individual depth level \\
No held-out validation & A fit can improve the same data used to create it and still be wrong & Reserve reference profiles or time blocks for evaluation \\
No proposal/review boundary & A numerical output could accidentally become an operational correction & Write a separate immutable proposal artefact; never edit official tables automatically \\
\bottomrule
\end{longtable}

\subsection{Out of scope}

The following are outside the scope of this framework:

\begin{itemize}[leftmargin=1.5em]
  \item fully automatic delayed-mode approval of corrections;
  \item changing or reproducing the upstream CORA reanalysis/validation system;
  \item redesigning the full MEOP processing and publication chain;
  \item treating expert historical coefficients as unquestionable truth;
  \item a sophisticated front-following ocean state estimate for every region;
  \item correction of pressure, oxygen, chlorophyll or light sensors.
\end{itemize}

\section{Existing and proposed data structures}

\subsection{MEOP target profiles}

Processed files follow

\begin{lstlisting}[style=meopcode]
data/data_prof/<deployment>/<smru_name>_<qf>_prof.nc
\end{lstlisting}

For example, an \pkg{hr1} file has dimensions \pkg{N_PROF} and \pkg{N_LEVELS}, with
one-dimensional profile metadata and two-dimensional hydrographic variables:

\begin{center}
\begin{tabularx}{0.96\textwidth}{>{\bfseries}p{0.25\textwidth}p{0.25\textwidth}X}
\toprule
Group & Typical variables & Use in this project \\
\midrule
Profile identity & \pkg{CYCLE_NUMBER}, platform metadata & Stable target-profile key \\
Time and position & \pkg{JULD}, \pkg{LATITUDE}, \pkg{LONGITUDE} plus QC & Match centre and filters \\
Raw hydrography & \pkg{PRES}, \pkg{TEMP}, \pkg{PSAL} plus QC & Primary data for bias estimation \\
Adjusted hydrography & \pkg{PRES_ADJUSTED}, \pkg{TEMP_ADJUSTED}, \pkg{PSAL_ADJUSTED} plus QC/error & Before/after evaluation only \\
Calibration metadata & Calibration-coefficient string array & Provenance and consistency checks \\
\bottomrule
\end{tabularx}
\end{center}

Use \pkg{hr1} when available because thermal-lag and high-resolution processing are
important upstream steps, but read the unadjusted \pkg{TEMP} and \pkg{PSAL} fields for
coefficient estimation. The processing level, choice of raw/adjusted field and all QC
rules must be stored in the output provenance.

\subsection{CORA reference profiles}

The CORA 202511 reference interface uses 540 populated 10-degree NetCDF tiles. The
preparation command uses one time origin and preserves the following fields for later
matchup work:

\begin{itemize}[leftmargin=1.5em]
  \item profile identifiers: \pkg{PLATFORM_NUMBER}, \pkg{CYCLE_NUMBER} and
  \pkg{DC_REFERENCE};
  \item source identifiers: \pkg{SOURCE_FILE}, \pkg{SOURCE_FILE_CLASS},
  \pkg{SOURCE_DATASET}, \pkg{REFERENCE_KIND} and \pkg{PROBE_TYPE};
  \item normalised \pkg{JULD}, \pkg{LATITUDE} and \pkg{LONGITUDE};
  \item \pkg{DATA_MODE}, \pkg{WMO_INST_TYPE}, and explicit
  \pkg{TEMP_SOURCE}/\pkg{PSAL_SOURCE}/\pkg{PRES_SOURCE} branch metadata;
  \item QC-masked, interpolated \pkg{TEMP} and \pkg{PSAL};
  \item the shared \pkg{PRES_GRID} coordinate.
\end{itemize}

An extended reader should return an \pkg{xarray.Dataset} or a small typed object that keeps
these fields together and decodes each tile's time before concatenation. A plain
dictionary of six arrays is no longer sufficient once matchups must be audited. Results
must record the CORA source download and preparation manifest hashes.

\subsection{Recommended comparison products}

Avoid making one enormous table with a row for every profile, reference neighbour and
pressure level. A two-part structure is easier to understand and more compact.

\subsubsection{Pair index: one row per target/reference pair}

Write a CSV for each tag or batch partition containing:

\begin{lstlisting}[style=meopcode]
target_tag, target_cycle, target_time, target_lat, target_lon,
reference_source, reference_platform, reference_cycle, reference_time,
reference_lat, reference_lon, distance_km, time_days,
space_weight, time_weight, combined_weight, accepted, rejection_reason
\end{lstlisting}

This table answers ``which observations influenced this comparison?'' without
duplicating vertical profiles. CSV is deliberately suggested because it is already
supported and can be inspected without a specialised stack. If scale becomes a problem,
Parquet can be considered later with an explicit dependency decision.

\subsubsection{Profile comparison NetCDF: one composite per target profile}

Write one NetCDF per tag with dimensions \pkg{N_PROF} and \pkg{N_LEVELS}. Suggested
variables are:

\begin{center}
\begin{tabularx}{0.96\textwidth}{>{\bfseries}p{0.31\textwidth}X}
\toprule
Variable/group & Meaning \\
\midrule
\pkg{PRES_GRID} & Common comparison pressure grid \\
\pkg{TARGET_TEMP_RAW}, \pkg{TARGET_PSAL_RAW} & QC-masked target values used for estimation \\
\pkg{TARGET_TEMP_ADJUSTED}, \pkg{TARGET_PSAL_ADJUSTED} & Optional operational values used only for evaluation \\
\pkg{REFERENCE_TEMP}, \pkg{REFERENCE_PSAL} & Weighted robust reference composite for this target profile \\
\pkg{REFERENCE_*_SPREAD} & Robust spread among reference neighbours, for example weighted MAD \\
\pkg{ANOMALY_*_RAW} & Raw target minus reference \\
\pkg{ANOMALY_*_ADJUSTED} & Adjusted target minus reference \\
\pkg{N_REFERENCE} & Number of accepted reference profiles for each target \\
\pkg{MEDIAN_DISTANCE_KM}, \pkg{MEDIAN_TIME_DAYS} & Matchup quality summaries \\
\pkg{MATCHUP_QC} & Good, marginal, no reference, bad target metadata, or insufficient overlap \\
\bottomrule
\end{tabularx}
\end{center}

Global attributes should include software version/commit, creation time, source file,
reference archive, variable policy, QC policy, pressure grid, matching thresholds,
weighting formula and a serialised configuration hash.

\subsubsection{Tag summary and coefficient proposal}

A batch-level summary CSV should contain coverage and robust anomalies by tag, variable,
pressure band and time period. A separate proposal CSV may contain:

\begin{lstlisting}[style=meopcode]
smru_platform_code, variable, model, T1, T2, S1, S2,
ci_low, ci_high, n_target_profiles, n_reference_profiles,
coverage_fraction, validation_delta_mad, sensitivity_range,
decision, decision_reason, method_version, created_at
\end{lstlisting}

The \pkg{decision} field should be one of \pkg{candidate}, \pkg{review},
\pkg{insufficient_data} or \pkg{reject}. This file is evidence for an operator; it is
not an alternative \pkg{table_coeff.csv}.

\subsection{Suggested output location}

Keep experimental products separate from standard plots until the method is approved:

\begin{lstlisting}[style=meopcode]
data/comparisons/cora/<method_version>/
  by_tags/<deployment>/<tag>_pairs.csv
  by_tags/<deployment>/<tag>_comparison.nc
  by_tags/<deployment>/<tag>_summary.csv
  runs/<timestamp>/summary.csv
  runs/<timestamp>/summary.md
  runs/<timestamp>/logs/<tag>.log
  latest/status.json
  proposals/coefficients_<timestamp>.csv
\end{lstlisting}

This mirrors the repository's existing resumable batch conventions without mixing
experimental coefficient suggestions into operational tables.

\section{General scientific methodology}

\subsection{Step 1: define a representative pilot set}

Begin with 6--10 tags, then expand to 30--50 before attempting the full archive. The
pilot must deliberately include:

\begin{itemize}[leftmargin=1.5em]
  \item zero-correction tags and tags with known constant corrections;
  \item at least one pressure-dependent coefficient;
  \item at least one time-varying salinity correction;
  \item contrasting regions, seasons, track lengths and profile counts;
  \item a missing-position or no-reference case to exercise failure reporting;
  \item one or more tags excluded or flagged by historical expert review.
\end{itemize}

Useful starting examples already present locally are shown below. They are chosen to
exercise code paths, not because historical coefficients are exact truth.

\begin{center}
\begin{tabularx}{0.98\textwidth}{>{\bfseries}p{0.24\textwidth}p{0.18\textwidth}X}
\toprule
Tag & Historical feature & Why useful \\
\midrule
\pkg{BM01-544-12} & All fixed coefficients zero & Negative-control case; 183 profiles \\
\pkg{ct88-225-12} & \pkg{S2=-0.4} & Strong constant salinity correction; existing regression fixture; 306 profiles \\
\pkg{ct62-01-10} & Non-zero T and S pressure terms & Exercises recovery of small slopes; 324 profiles \\
\pkg{ct84-533-12} & Fixed coefficients plus time-varying salinity & Exercises separation of a tag-wide model from profile-index change; 321 profiles \\
\pkg{ct153-Pendragon-19} & Non-zero \pkg{T2} and \pkg{S2} & Long 1,093-profile track; already an LR-only adjustment regression case \\
\pkg{ct104-EF970-13} & Small constant T/S corrections & Long Southern Ocean track with 1,171 profiles \\
\bottomrule
\end{tabularx}
\end{center}

Confirm the final pilot list after producing a concise inventory of reference coverage
and historical decisions.

\subsection{Step 2: choose explicit input policies}

The comparison must make three choices explicit:

\begin{enumerate}[leftmargin=1.6em]
  \item \textbf{Target processing level.} Start with \pkg{hr1} where present, with a
  documented fallback. Do not silently mix levels in one scientific result.
  \item \textbf{Target variable branch.} Use raw \pkg{TEMP}/\pkg{PSAL} for estimating
  calibration bias. Carry adjusted variables separately to measure whether existing
  corrections improve external agreement.
  \item \textbf{Reference variable branch.} Use the prepared CORA best-parameter policy:
  adjusted whole-profile values when available, otherwise raw. Preserve
  \pkg{TEMP_SOURCE}, \pkg{PSAL_SOURCE}, \pkg{PRES_SOURCE} and \pkg{DATA_MODE} in every
  auditable result. A different branch policy requires a separately identified reference
  build and sensitivity comparison.
\end{enumerate}

\subsection{Step 3: quality control and independence}

Decode byte, character and object QC arrays through one tested helper. At minimum:

\begin{itemize}[leftmargin=1.5em]
  \item reject invalid target time or position and apply target T/S QC level by level;
  \item verify the reference manifest and enforce the prepared archive's accepted QC
  policy, then reject any remaining non-finite reference values;
  \item reject missing/fill values and non-physical pressure ordering;
  \item require a minimum common pressure span and number of valid comparison levels;
  \item retain the reason for every rejection;
  \item exclude exact duplicates and verify that no MEOP observation is compared with
  itself through the reference archive.
\end{itemize}

The local \pkg{list_wmo_numbers.csv} and CORA fields such as \pkg{PLATFORM_NUMBER} and
\pkg{DC_REFERENCE} can help with independence checks. If complete self-source exclusion
cannot be guaranteed, that limitation must be visible in the result and discussed with
the scientific reviewer before coefficient estimation.

\subsection{Step 4: create space--time neighbours}

Convert longitudes to a consistent convention and positions to unit-sphere Cartesian
coordinates. A \pkg{scipy.spatial.cKDTree} search on the sphere, followed by an exact
haversine distance, is efficient and naturally handles the anti-meridian. Time should
be normalised to UTC before matching.

Reasonable \emph{initial values for sensitivity testing}, not final scientific constants,
are:

\begin{itemize}[leftmargin=1.5em]
  \item maximum distance: 150 km;
  \item maximum absolute time difference: 30 days;
  \item maximum neighbours retained per target: 20;
  \item minimum accepted reference profiles per target: 3.
\end{itemize}

Repeat key results for 50/100/200 km and 10/30/60 days. A coefficient
that changes substantially across these defensible settings should be marked for review,
not reported as a confident adjustment.

\subsection{Step 5: vertical interpolation and the reference composite}

Use a documented common pressure grid and never extrapolate outside the observed range
of a profile. Perform interpolation only after masking invalid values. Preserve a count
and robust spread at every target profile and pressure level.

A simple starting weight for accepted reference neighbour \(j\) is

\begin{equation}
w_j = \exp\left[-\frac{1}{2}\left(\frac{d_j}{L}\right)^2
                 -\frac{1}{2}\left(\frac{\Delta t_j}{\tau}\right)^2\right],
\end{equation}

where \(d_j\) is distance, \(\Delta t_j\) is time separation, and \(L\), \(\tau\) are
configurable scales. Compare a weighted mean with a weighted median or trimmed mean.
The robust choice is preferable where enough profiles exist; the method should not
pretend that a one-neighbour composite has well-determined reference uncertainty.

\subsection{Step 6: compute and inspect anomalies}

For each target profile \(i\), pressure \(p\), and variable \(X\), compute

\begin{equation}
\Delta X_i(p)=X_{\mathrm{MEOP},i}(p)-\widehat{X}_{\mathrm{ref},i}(p).
\end{equation}

Summaries should retain profile structure. Useful diagnostics include:

\begin{itemize}[leftmargin=1.5em]
  \item anomaly versus pressure, with median and interquartile/reference-spread bands;
  \item profile-median anomaly versus time/profile number;
  \item anomaly versus distance, time separation and number of neighbours;
  \item maps coloured by anomaly and matchup coverage;
  \item raw-versus-adjusted residual comparison;
  \item summaries by stable pressure bands and, where scientifically useful, by
  potential-density or water-mass class;
  \item results under each space/time sensitivity setting.
\end{itemize}

First explain visible structure as oceanography before labelling it
sensor bias. Front crossings, seasonal cycles, shallow mixed layers, sea-ice regions and
sparse reference coverage are all plausible causes of structured anomalies.

\subsection{Step 7: optional robust coefficient estimation}

Only after the comparison passes an explicit scientific review, fit the existing
operational parameterisation to raw
anomalies over agreed stable pressure ranges:

\begin{align}
\Delta T_i(p) &= 10^{-3}T1\,p + T2 + \epsilon_{i,p}, \\
\Delta S_i(p) &= 10^{-3}S1\,p + S2 + \epsilon_{i,p}.
\end{align}

Use robust regression so a small number of water-mass mismatches cannot dominate the
fit. \pkg{scipy.optimize.least_squares} with a robust loss, or a clearly implemented
iteratively reweighted method, avoids adding a large machine-learning dependency.

Do not treat all depth samples as independent. Estimate confidence by resampling target
profiles or multi-day blocks. Store the full coefficient distribution or at least its
percentile interval. Use separate training and validation profiles or time blocks.

For time-varying salinity, first plot profile-level residuals after the fixed model. A
piecewise constant/linear correction with very few change points may be investigated as
a stretch task. Penalise extra breakpoints, require persistence, bootstrap breakpoint
stability and respect the existing limit of four index/offset pairs if proposing an
operationally compatible result.

\subsection{Step 8: decision gates}

A proposed coefficient should be labelled \pkg{candidate} only when all agreed gates
pass. Starting gates to refine with the scientific review group are:

\begin{itemize}[leftmargin=1.5em]
  \item enough matched target profiles and adequate deployment-time coverage;
  \item more than one reference platform and reasonable spatial/time separations;
  \item sufficient vertical overlap in the chosen pressure band;
  \item a bootstrap interval narrow enough for the intended correction;
  \item the sign and magnitude are stable across matching sensitivities;
  \item held-out median bias and MAD improve, not just training fit;
  \item no implausible T/S or density behaviour is introduced;
  \item results are understandable on the diagnostic plots.
\end{itemize}

Failure of a gate is useful information. The correct output may be
\pkg{insufficient_data} or \pkg{review}; the estimator must never force every tag to have
a non-zero correction.

\section{Validation plan}

\subsection{Synthetic recovery tests}

Create small synthetic MEOP/CORA fixtures with known offsets, pressure slopes, missing
levels, bad QC, time offsets, longitude wrap and a known change point. Tests should show
that the pipeline:

\begin{itemize}[leftmargin=1.5em]
  \item selects exactly the expected neighbours;
  \item handles 179 degrees east versus 179 degrees west correctly;
  \item never extrapolates vertically;
  \item recovers injected \pkg{T1/T2/S1/S2} within a stated tolerance;
  \item does not infer a correction in the zero-bias case;
  \item reports no-reference and insufficient-data cases honestly;
  \item widens uncertainty when reference coverage is poor.
\end{itemize}

\subsection{Historical comparison}

Stratify historical tags into zero coefficient, constant offset, pressure-dependent,
time-varying, salinity removed and fully removed groups. For each stratum, compare
candidate results with the operator tables, while recognising that disagreement can mean
the new method is wrong, the reference is unsuitable, or the historical decision used
evidence not available to the algorithm.

Report at least coefficient sign agreement, robust difference, interval coverage,
false-positive rate on zero-coefficient tags, and the fraction classified as
insufficient. Review individual large disagreements rather than hiding them in one
correlation coefficient.

\subsection{Held-out and sensitivity validation}

Fit on one part of the deployment and evaluate on another, or divide independent
reference platforms between fit and evaluation. Compare before/after median anomaly and
MAD. Repeat under several matching windows and reference-variable policies. The most
important result is not the smallest residual; it is a correction that generalises and
remains stable when defensible assumptions change.

\subsection{Scientific review gate}

Before coefficient estimation begins, review:

\begin{enumerate}[leftmargin=1.6em]
  \item corrected CORA coverage and explicit batch status counts;
  \item pair maps and anomaly plots for the pilot tags;
  \item results of matching-window sensitivity tests;
  \item evidence that raw and adjusted branches are not mixed;
  \item a recommendation to proceed with fixed coefficient estimation, limit the
  project to comparison, or revise the reference strategy.
\end{enumerate}

Proceeding to estimation should be a scientific decision made at this review, not an
automatic consequence of completing the comparison software.

\section{Final perspective}

Most of the engineering substrate is already present: processed MEOP profiles,
coefficient application, CORA tiles, plotting, batch state and tests. The missing piece
is the scientific middle layer that turns ``profiles in the same broad box'' into
traceable, independent and uncertainty-aware comparisons.

That layer is a coherent scientific development project spanning data handling,
geospatial algorithms, oceanographic interpretation, statistics, testing and
reproducible software development. The systematic comparison is useful in its own
right and provides the evidence needed to decide how far defensible adjustment
estimation can go.

\begin{thebibliography}{9}

\bibitem{roquet2011}
Roquet, F., Charrassin, J.-B., Marchand, S., Boehme, L., Fedak, M., Reverdin, G., and
Guinet, C. (2011).
Validating hydrographic data obtained from seal-borne satellite-relayed data loggers.
\emph{Journal of Atmospheric and Oceanic Technology}, 28, 787--801.
\href{https://doi.org/10.1175/2010JTECHO801.1}{doi:10.1175/2010JTECHO801.1}.

\bibitem{meopmanual}
MEOP/Sea Mammal data-processing manual, version 1.2, available locally as
\pkg{docs/seamammal_user_manual_version1.2.pdf}.

\bibitem{toolboxstatus}
MEOP Process renovation status and toolbox guide, available locally as
\pkg{docs/meop_toolbox_status.tex}.

\bibitem{corapum}
Copernicus Marine Service, Product User Manual for
\pkg{INSITU_GLO_PHY_TS_DISCRETE_MY_013_001}, issue 1.17 (November 2025).
\href{https://documentation.marine.copernicus.eu/PUM/CMEMS-INS-PUM-013-001.pdf}{Official CORA user manual}.

\bibitem{coradata}
Copernicus Marine Service, Global Ocean CORA in-situ observations, delayed mode.
\href{https://data.marine.copernicus.eu/product/INSITU_GLO_PHY_TS_DISCRETE_MY_013_001/services}{Product and data access page}.

\end{thebibliography}

\end{document}
